\documentclass[11pt,a4paper]{article}
\usepackage[margin=2.3cm]{geometry}
\usepackage{booktabs}
\usepackage{graphicx}
\usepackage[hidelinks]{hyperref}
\usepackage{microtype}
\usepackage{siunitx}
\usepackage{float}
\usepackage[section]{placeins}

\IfFileExists{generated_metrics.tex}{
  \input{generated_metrics.tex}
}{
  \PackageError{ai-art-report}
    {Measured result fragments are missing}
    {Run the experiments and scripts/make_latex_table.py before compiling.}
}

\title{Binary Detection of AI-Generated Art Using Transfer Learning and
Paper-Inspired ConvNeXt Attention}
\author{Your Name}
\date{\today}

\begin{document}
\maketitle

\begin{abstract}
This project evaluates binary classification of human and AI-generated artwork
using AI-ArtBench. Five ImageNet-pretrained configurations compare MobileNetV2
and ConvNeXt-Tiny under augmentation, partial fine-tuning, and a lightweight
squeeze-and-excitation ablation. The validation-selected ConvNeXt-Tiny achieved
test accuracy \BestTestAccuracy, F1 \BestTestFOne, and ROC-AUC \BestTestAUC.
Robustness, subgroup errors, and Grad-CAM qualify these results.
\end{abstract}

\section{Introduction}
Diffusion models can produce artwork that is difficult to distinguish from
human-created work. Silva et al.\ introduced AI-ArtBench and
AttentionConvNeXt for style and source classification
\cite{silva2024artbrain}. This project narrows the target to the binary
question: human-created or AI-generated? The goal is a reproducible comparison
of practical transfer-learning models, with explicit checks for source and
image-processing shortcuts.

\section{Dataset and protocol}
AI-ArtBench \cite{silva2023aiartbench} combines 60,000 human artworks from ArtBench-10
\cite{liao2022artbench} with 125,015 images generated by Latent Diffusion and
Stable Diffusion. The published partition contains 155,015 training and 30,000
test images. Human and Latent Diffusion images are \(256\times256\), whereas
Stable Diffusion images are \(768\times768\).

The published test partition is kept isolated. A deterministic seed-42 sample
contains \TrainCount{} training, \ValidationCount{} validation, and
\TestCount{} test images. Every partition is binary-balanced, each AI half is
balanced between the two generators, and all ten styles receive equal quotas.
Validation is drawn only from the published training partition. No test image
is used for optimization, early stopping, or model selection.

\section{Methods}
Images are converted to RGB, resized to \(224\times224\), and normalized with
ImageNet statistics. Training augmentation uses random resized crops,
horizontal flips, mild colour jitter, and occasional autocontrast.

MobileNetV2 \cite{sandler2018mobilenetv2} is the efficient baseline and
ConvNeXt-Tiny \cite{liu2022convnext} is the larger backbone. Frozen-backbone
and last-stage fine-tuning settings are compared. The final ablation applies
SE channel weighting \cite{hu2018squeeze} to the last ConvNeXt feature block.
It is a paper-inspired hypothesis test, not a reproduction of the full
multi-level AttentionConvNeXt architecture.

All models use a single-logit binary head, binary cross-entropy, AdamW, batch
size 32, and at most eight epochs. The learning rates are \(10^{-3}\) for
head-only training and \(10^{-4}\) for partial fine-tuning. Training stops
after three validation epochs without an F1 improvement. The classification
threshold is fixed at 0.5.

\section{Results}
\begin{table}[ht]
\centering
\resizebox{\textwidth}{!}{\input{generated_results_table.tex}}
\caption{Validation-model-selection and held-out test results.}
\label{tab:results}
\end{table}

The test-F1 changes associated with augmentation, changing the frozen backbone,
partial fine-tuning, and the SE ablation were respectively
\AugmentationDelta, \BackboneDelta, \FinetuningDelta, and \AttentionDelta.
These controlled comparisons are more informative than treating the largest
model as automatically preferable.

\begin{figure}[ht]
\centering
\includegraphics[width=.47\textwidth]{generated_confusion.png}
\hfill
\includegraphics[width=.47\textwidth]{generated_roc.png}
\caption{Confusion matrix and ROC curve for the validation-selected model.}
\end{figure}

\IfFileExists{generated_replication_table.tex}{
\section{Independent replication and overfitting audit}
After training and primary test reporting were complete, a second
\ReplicationCount-image holdout was sampled from the unused official test
pool with seed 4242. It has the same source/style quotas as the first test set
and zero overlap with training, validation, or primary-test paths. All five
saved checkpoints were evaluated without retraining or threshold changes.
Both clean-training and replication measurements use only direct
\(224\times224\) resizing and ImageNet normalization; the 128-pixel robustness
resampling path is not used in this audit.

\begin{table}[ht]
\centering
\resizebox{\textwidth}{!}{\input{generated_replication_table.tex}}
\caption{Clean-training, primary-test, and independent-replication F1.
Intervals use 1,000 source/style-stratified bootstrap resamples.}
\label{tab:replication}
\end{table}

The validation-selected E3 model obtained replication F1
\SelectedReplicationFOne{} (95\% CI
\([\SelectedReplicationFOneCILow,\SelectedReplicationFOneCIHigh]\)).
Its clean-training minus replication gap was
\SelectedTrainReplicationGap{} (95\% CI
\([\SelectedGapCILow,\SelectedGapCIHigh]\)), while replication minus the
primary test changed F1 by \SelectedReplicationOriginalDelta. These values
measure same-dataset overfitting and sampling stability, not transfer to an
unseen generator. The highest replication F1 belonged to
\ReplicationBestExperiment.
}{}

\section{Subgroup and robustness analysis}
Source groups contain only one binary class, so source-level F1 would be
misleading. Table~\ref{tab:source} instead reports false-positive rate for
Human and false-negative rate for each AI generator. The highest source error
was \HighestSourceError{} for \HighestErrorSource.

\begin{table}[ht]
\centering
\input{generated_source_table.tex}
\caption{Class-appropriate test error by image source.}
\label{tab:source}
\end{table}

Each style contains both binary classes and therefore supports the full metric
set. F1 ranged from \LowestStyleFOne{} for \LowestStyle{} to
\HighestStyleFOne{} for \HighestStyle.

\begin{table}[ht]
\centering
\resizebox{\textwidth}{!}{\input{generated_style_table.tex}}
\caption{Held-out test metrics by artistic style.}
\end{table}

Contrast, JPEG compression, and a common 128-pixel resampling bottleneck test
sensitivity to post-processing and resolution artefacts. The worst tested
condition was \WorstRobustnessCondition{} at value
\WorstRobustnessValue, changing F1 by \WorstRobustnessDelta{} relative to the
clean test set.

\begin{table}[ht]
\centering
\resizebox{\textwidth}{!}{\input{generated_robustness_table.tex}}
\caption{Robustness results for the validation-selected model.}
\end{table}

\begin{figure}[H]
\centering
\includegraphics[width=.60\textwidth]{generated_robustness.png}
\caption{F1 under clean and perturbed test conditions.}
\end{figure}

\section{Explainability}
Grad-CAM \cite{selvaraju2017gradcam} visualizes spatial sensitivity for
representative correct and incorrect predictions. These heatmaps can reveal
attention to texture, borders, signatures, or compression patterns, but they
are not causal explanations.

\IfFileExists{generated_gradcam_misclassified.png}{
\begin{figure}[H]
\centering
\includegraphics[width=.68\textwidth]{generated_gradcam_misclassified.png}
\caption{Grad-CAM examples for misclassified test images.}
\end{figure}
}{}

\section{Limitations}
AI-ArtBench contains only two related diffusion families and may reward
generator-specific artefacts rather than a general concept of artificiality.
The human images are digitized artworks, while generated images are
born-digital, and source resolutions differ. Balanced subsampling improves
interpretability but the primary protocol uses only 10,000 of 185,015 images;
the separate audit adds 2,000 same-distribution test images. Generalization to
unseen generators, edited images, photographs, and new human-art collections
is not established. Grad-CAM should not be interpreted as proof of model
reasoning.

\section{Conclusion}
The validation-selected model achieved held-out F1 \BestTestFOne{} and
ROC-AUC \BestTestAUC{} under the fixed official-split protocol. The ablations
show how augmentation, backbone choice, fine-tuning, and SE weighting changed
performance, while the subgroup and perturbation results delimit the strength
of the headline score. The resulting detector is a controlled coursework
benchmark rather than a universal artwork-authentication system.
\IfFileExists{generated_replication_table.tex}{
The disjoint replication holdout provides a separate estimate of
same-distribution generalization without changing the original model
selection or headline result.
}{}

\bibliographystyle{plain}
\bibliography{references}
\end{document}
